\clearpage\pagestyle{otMain}\setlength{\otBlocEspacement}{2.5mm}\fontsize{11}{15}\selectfont\corrigesChapitre{15}{Symétries et demi-tour}
\begin{blocLivre}{cor-15-1}\begin{corrige}{1}{Un axe vertical}\demarche{Reporter la distance perpendiculaire à l’axe.}
Compter trois carreaux entre A et l’axe. Reporter trois carreaux de l’autre côté sur la même horizontale. A′ a la même ordonnée que A.
\label{fig-cor-15-1}\figureGeo[corrige]{sym-1}
\end{corrige}\end{blocLivre}
\begin{blocLivre}{cor-15-2}\begin{corrige}{2}{Un axe horizontal}\demarche{Garder la direction perpendiculaire à l’axe.}
B est deux carreaux au-dessus. B′ est deux carreaux au-dessous, avec la même abscisse.
\label{fig-cor-15-2}\figureGeo[corrige]{sym-2}
\end{corrige}\end{blocLivre}
\begin{blocLivre}{cor-15-3}\begin{corrige}{3}{Un demi-tour}\demarche{Faire du centre le milieu.}
Prolonger (AO) au-delà de O et reporter OA. Le point obtenu vérifie $OA=OA′$ et O appartient à [AA′].
\label{fig-cor-15-3}\figureGeo[corrige]{sym-3}
\end{corrige}\end{blocLivre}
\begin{blocLivre}{cor-15-4}\begin{corrige}{4}{Le centre lui-même}\demarche{Appliquer la définition au centre.}
Le point O reste à sa place : son image est $\res O$. Sa distance au centre vaut zéro.
\label{fig-cor-15-4}\figureGeo[corrige]{sym-4}
\end{corrige}\end{blocLivre}
\begin{blocLivre}{cor-15-5}\begin{corrige}{5}{Longueur conservée}\demarche{Utiliser la conservation.}
La symétrie centrale conserve les longueurs : $\res{A′B′=4{,}5\ \mathrm{cm}}$.
\label{fig-cor-15-5}\figureGeo[corrige]{sym-5}
\end{corrige}\end{blocLivre}
\begin{blocLivre}{cor-15-6}\begin{corrige}{6}{Axe diagonal}\demarche{Contrôler le milieu et la perpendicularité.}
La diagonale échange les déplacements horizontaux et verticaux. L’image est $\res{A′(3;1)}$. Le milieu de [AA′] est $(2;2)$ sur l’axe et [AA′] lui est perpendiculaire.
\label{fig-cor-15-6}\figureGeo[corrige]{sym-6}
\end{corrige}\end{blocLivre}
\begin{blocLivre}{cor-15-7}\begin{corrige}{7}{Image d’un triangle}\demarche{Transformer les trois sommets.}
Construire A′, B′ et C′ en faisant de O le milieu des trois segments correspondants, puis relier les images. Le triangle image est de même forme et de même taille.
\label{fig-cor-15-7}\figureGeo[corrige]{sym-7}
\end{corrige}\end{blocLivre}
\begin{blocLivre}{cor-15-8}\begin{corrige}{8}{Retrouver le centre}\demarche{Inverser la définition.}
Le centre est le $\res{\text{milieu de [AA′]}}$. On peut le construire à la règle et au compas avec la médiatrice du segment.
\label{fig-cor-15-8}\figureGeo[corrige]{sym-8}
\end{corrige}\end{blocLivre}
\begin{blocLivre}{cor-15-9}\begin{corrige}{9}{Image d’une droite}\demarche{Utiliser deux points pour déterminer une droite.}
Construire A′ et B′ puis tracer (A′B′). L’image d’une droite par un demi-tour est $\res{\text{parallèle}}$ à la droite de départ.
\label{fig-cor-15-9}\figureGeo[corrige]{sym-9}
\end{corrige}\end{blocLivre}
\begin{blocLivre}{cor-15-10}\begin{corrige}{10}{Un cercle image}\demarche{Transformer le centre et garder le rayon.}
Construire C′, image de C, puis le cercle de centre C′ et de rayon $\res{2\ \mathrm{cm}}$. Le rayon est conservé.
\label{fig-cor-15-10}\figureGeo[corrige]{sym-10}
\end{corrige}\end{blocLivre}
\begin{blocLivre}{cor-15-11}\begin{corrige}{11}{Une construction sur feuille blanche}\demarche{Aligner les points et reporter une longueur.}
Tracer (AO). Prendre au compas l’ouverture OA et la reporter depuis O sur la demi-droite opposée à [OA). L’autre intersection avec le cercle de centre O et rayon OA est A′.
\label{fig-cor-15-11}\figureGeo[corrige]{sym-11}
\end{corrige}\end{blocLivre}
\begin{blocLivre}{cor-15-12}\begin{corrige}{12}{Aire et périmètre}\demarche{Distinguer déplacement et changement d’échelle.}
Les longueurs sont conservées, donc le périmètre reste $\res{18\ \mathrm{cm}}$. L’aire est aussi conservée : $\res{12\ \mathrm{cm}^2}$. La position a changé, pas la taille.
\label{fig-cor-15-12}\figureGeo[corrige]{sym-12}
\end{corrige}\end{blocLivre}
\begin{blocLivre}{cor-15-13}\begin{corrige}{13}{Centre d’un rectangle}\demarche{Chercher des milieux communs.}
Les diagonales ont le même milieu O. Le demi-tour de centre O échange $\res{A\text{ et }C}$, $\res{B\text{ et }D}$. Les côtés correspondants se superposent ; le rectangle est sa propre image.
\label{fig-cor-15-13}\figureGeo[corrige]{sym-13}
\end{corrige}\end{blocLivre}
\begin{blocLivre}{cor-15-14}\begin{corrige}{14}{Deux demi-tours}\demarche{Appliquer deux fois la définition du milieu.}
Si A devient A′, O est le milieu de [AA′]. En repartant de A′, le même demi-tour renvoie donc sur A. Cela vaut pour tous les sommets : on retrouve $\res{\text{le triangle initial}}$.
\label{fig-cor-15-14}\figureGeo[corrige]{sym-14}
\end{corrige}\end{blocLivre}
\begin{blocLivre}{cor-15-15}\begin{corrige}{15}{Un centre commun à vérifier}\demarche{Comparer les milieux des segments d’images.}
Le milieu de [AA′], avec $A(-3;1)$ et $A′(3;-1)$, est le point $(0;0)$. Le milieu de [BB′], avec $B(-1;3)$ et $B′(1;-1)$, est le point $(0;1)$. Ils diffèrent : $\res{\text{aucune même symétrie centrale}}$ ne réalise les deux échanges, car son centre devrait être le milieu des deux segments.
\label{fig-cor-15-15}\figureGeo[corrige]{sym-15}
\end{corrige}\end{blocLivre}
